\documentclass[12pt,a4paper]{article}
\usepackage[utf8]{inputenc}
\usepackage[brazilian]{babel}
\usepackage{amsmath}
\usepackage{amsfonts}
\usepackage{amssymb}
\usepackage{graphicx}
\usepackage{hyperref}
\usepackage{geometry}
\usepackage{listings}
\usepackage{booktabs}
\usepackage{xcolor}
\usepackage{setspace} % Controle de espaçamento entre linhas
\usepackage{indentfirst} % Indenta o primeiro parágrafo de cada seção
\usepackage{times} % Fonte Times New Roman

% Configurações de margens da ABNT: Esquerda e Superior (3cm), Direita e Inferior (2cm)
\geometry{
    a4paper,
    top=3.0cm,
    bottom=2.0cm,
    left=3.0cm,
    right=2.0cm
}

% Espaçamento de 1.5 cm conforme norma ABNT
\onehalfspacing

% Recuo do parágrafo padrão ABNT (1,25 cm)
\setlength{\parindent}{1.25cm}

% Estilo para trechos de código
\definecolor{codegreen}{rgb}{0,0.6,0}
\definecolor{codegray}{rgb}{0.5,0.5,0.5}
\definecolor{codepurple}{rgb}{0.58,0,0.82}
\definecolor{backcolour}{rgb}{0.95,0.95,0.92}

\lstdefinestyle{mystyle}{
    backgroundcolor=\color{backcolour},   
    commentstyle=\color{codegreen},
    keywordstyle=\color{blue},
    numberstyle=\tiny\color{codegray},
    stringstyle=\color{codepurple},
    basicstyle=\ttfamily\footnotesize,
    breakatwhitespace=false,         
    breaklines=true,                 
    captionpos=b,                    
    keepspaces=true,                 
    numbers=left,                    
    numbersep=5pt,                  
    showspaces=false,                
    showstringspaces=false,
    showtabs=false,                  
    tabsize=2
}
\lstset{style=mystyle}

\begin{document}

% --- CAPA CONFORME ABNT ---
\begin{titlepage}
    \begin{center}
        {\large\bfseries UNIVERSIDADE DO ESTADO DA BAHIA -- UNEB} \\
        {\large DEPARTAMENTO DE CIÊNCIAS EXATAS E DA TERRA} \\
        {\large CURSO DE ENGENHARIA DE COMPUTAÇÃO / SISTEMAS DE INFORMAÇÃO} \\
        \vspace{4.0cm}
        {\large\bfseries GRUPO 5} \\
        \vspace{4.0cm}
        {\Large\bfseries ASSISTENTE DE CONSULTA FUNDAMENTADA SOBRE O DGP DO CNPQ: UMA ABORDAGEM UTILIZANDO RAG, EMBEDDINGS VETORIAIS E VALIDAÇÃO DE SUFICIÊNCIA} \\
        \vspace{4.0cm}
        \vfill
        {\large Salvador -- BA} \\
        {\large 2026}
    \end{center}
\end{titlepage}

% --- FOLHA DE ROSTO ---
\begin{titlepage}
    \begin{center}
        {\large\bfseries GRUPO 5} \\
        \vspace{4.0cm}
        {\Large\bfseries ASSISTENTE DE CONSULTA FUNDAMENTADA SOBRE O DGP DO CNPQ: UMA ABORDAGEM UTILIZANDO RAG, EMBEDDINGS VETORIAIS E VALIDAÇÃO DE SUFICIÊNCIA} \\
        \vspace{3.0cm}
        
        \hspace{7.0cm}
        \begin{minipage}{8.0cm}
            \setstretch{1.0}
            \small Relatório técnico científico apresentado à disciplina de Inteligência Artificial do Departamento de Ciências Exatas e da Terra da Universidade do Estado da Bahia como requisito parcial para avaliação no Semestre 2026.1.\\
            \\
            Professor Orientador: Dr. Roberto Abreu.
        \end{minipage}
        
        \vfill
        {\large Salvador -- BA} \\
        {\large 2026}
    \end{center}
\end{titlepage}

\newpage
\tableofcontents
\newpage

\section{Definição do Problema}
O advento dos Modelos de Linguagem de Grande Escala (LLMs), tais como GPT e Llama, revolucionou o processamento de linguagem natural. Contudo, em cenários corporativos ou acadêmicos que demandam alta confiabilidade informacional, a utilização isolada de LLMs apresenta vulnerabilidades críticas. Dentre estas fragilidades, destaca-se a \textit{alucinação}, fenômeno pelo qual o modelo gera fatos incorretos ou inventados com alta convicção gramatical. Adicionalmente, estes modelos não possuem acesso em tempo real a bases privadas, dinâmicas ou de escopo restrito.

No contexto das instituições acadêmicas brasileiras, o Diretório de Grupos de Pesquisa (DGP) do CNPq e os perfis da Plataforma Lattes reúnem toda a produção intelectual nacional. No entanto, a consulta a estes portais é feita por buscadores convencionais que não permitem cruzamentos complexos nem diálogos inteligentes. O problema central consiste em desenvolver uma ferramenta que atue como assistente de busca inteligente capaz de responder a consultas complexas sobre a produção de pesquisadores e atividades dos grupos de forma estritamente fundamentada, evitando respostas falsas e atestando as fontes originais dos dados bibliográficos coletados.

\section{Base Documental}
A base de dados é mantida e populada localmente com uma fração amostral selecionada de dados brutos de grupos de pesquisa e pesquisadores de diversas instituições de ensino superior do estado da Bahia (tais como UESC, UFBA, UNEB, UNIFACS, etc.), extraídos diretamente das plataformas públicas do CNPq por um crawler customizado. A volumetria atende aos seguintes parâmetros definidos na especificação:
\begin{itemize}
    \item \textbf{Grupos de Pesquisa:} Armazena dados de espelhos completos de grupos de pesquisa contendo informações de ano de formação, repercussão, área geral, líderes e instituição.
    \item \textbf{Recursos Humanos:} Armazena a listagem de líderes, pesquisadores, técnicos e estudantes envolvidos.
    \item \textbf{Linhas de Pesquisa:} Objetivos, palavras-chave de indexação e setores de aplicação.
    \item \textbf{Currículos Lattes:} Extrai o resumo profissional do pesquisador, seus artigos publicados (incluindo ISSN, volume, páginas e ano de publicação), livros, capítulos de livros publicados e projetos de pesquisa.
\end{itemize}

Esta base é convertida em arquivos JSON brutos e ingerida para tabelas do banco de dados relacional PostgreSQL, totalizando mais de 100 documentos detalhados, o que ultrapassa os 80 chunks exigidos no edital.

\section{Pipeline de Ingestão e Processamento}
A ingestão e fragmentação dos dados são críticas para garantir que o contexto recuperado seja coeso e focado. A arquitetura de ingestão é composta por:
\begin{enumerate}
    \item \textbf{Watcher e Extrator de Eventos:} Um script baseado na biblioteca Python \texttt{watchdog} monitora o diretório de dados brutos (\texttt{data/raw-data}) em tempo real. Quando um arquivo JSON de grupo ou pesquisador é salvo, o ETL é disparado.
    \item \textbf{Normalização:} O ETL valida a integridade do JSON, converte campos numéricos (como anos de formação e páginas), limpa tags HTML das strings e cruza o ISSN dos artigos científicos com a base nacional do Qualis CAPES.
    \item \textbf{Sincronização de RAG:} O script \texttt{langchain\_api/ingest\_db.py} avalia quais entidades sofreram alteração desde a última vetorização (sincronização incremental).
    \item \textbf{Chunking e Segmentação:} O conteúdo textual de currículos Lattes e perfis de grupos é segmentado linearmente utilizando a estratégia de chunking baseada em caracteres. 
    \item \textbf{Parâmetros de Overlap e Justificativa:} O texto é dividido em chunks de no máximo 1000 caracteres, com uma sobreposição (overlap) de 200 caracteres. Essa escolha faz sentido para o problema porque 1000 caracteres comportam confortavelmente um parágrafo denso de informações cadastrais ou a listagem detalhada de uma produção acadêmica. O overlap de 200 caracteres é essencial para evitar o corte abrupto de palavras e garantir que entidades na fronteira do fragmento sejam preservadas em seu contexto semântico no chunk subsequente.
\end{enumerate}

\section{Arquitetura da Solução}
A arquitetura adota o estilo de monólito em Python para reduzir a latência de comunicação de rede. A estrutura do projeto é a seguinte:
\begin{itemize}
    \item \textbf{Camada de Scraping:} Sequencial e focada em uma única instância ativa de Playwright e BeautifulSoup para ler dados cadastrais, popups de integrantes e linhas de pesquisa sem disparar alertas de limite de requisições.
    \item \textbf{Camada de Persistência Relacional e Vetorial:} PostgreSQL com extensão \texttt{pgvector}. As tabelas relacionais gerenciam a integridade dos dados e a tabela \texttt{rag\_chunk} persiste os embeddings vetorizados na coluna vetorial de dimensão 1536.
    \item \textbf{Camada de Recuperação e Geração (FastAPI):} Expõe uma interface API REST com o endpoint \texttt{/question}.
\end{itemize}

O processamento da busca vetorial envolve os seguintes eixos operacionais:
\begin{itemize}
    \item \textbf{Vetorização da Consulta:} A consulta em linguagem natural feita pelo usuário é convertida em tempo de execução em um vetor denso de 1536 dimensões utilizando o modelo \texttt{text-embedding-3-small} da OpenAI.
    \item \textbf{Uso de Similaridade Vetorial:} A similaridade entre a consulta e os chunks no banco é mensurada através da menor distância Euclidiana (L2), representada pelo operador matemático \texttt{<->} no pgvector. Chunks menores na distância Euclidiana indicam maior proximidade semântica.
    \item \textbf{Número de Trechos Recuperados:} O banco de dados recupera e retorna os 5 trechos (\texttt{LIMIT 5}) mais similares à consulta.
    \item \textbf{Incorporação ao Prompt:} Os 5 chunks retornados são limpos e concatenados em uma string de contexto. Esse contexto é então injetado por meio de interpolação no prompt do sistema sob a variável \texttt{\{context\}}, delimitando estritamente as fronteiras da resposta do LLM.
\end{itemize}

\subsection{Estratégias de Validação e Mecanismo Híbrido}
Para garantir a confiabilidade das respostas e atender às diretrizes do edital, a solução incorpora as seguintes estratégias de validação e processamento híbrido:
\begin{itemize}
    \item \textbf{Recusa Explícita por Insuficiência de Contexto:} O prompt do sistema (\texttt{SIMPLE\_PROMPT}) impõe uma restrição comportamental estrita sobre o LLM. Caso o contexto vetorial injetado não contenha os fatos necessários para responder à pergunta (como no caso de questionamentos sobre astrofísica ou Albert Einstein), o modelo é instruído a recusar a resposta de forma explícita e padronizada (ex: \textit{"Não encontrei informações..."}), eliminando palpites ou alucinações.
    \item \textbf{Indicação de Fontes e Trechos Utilizados:} O fluxo de retorno da API não se limita ao texto gerado. A resposta do endpoint \texttt{/question} é encapsulada em um objeto de resposta que anexa um array contendo os metadados das fontes (\texttt{sources}) que compuseram o contexto (título do documento, tipo do documento Lattes/DGP, ID do registro e similaridade percentual). Isso garante a rastreabilidade e auditabilidade do sistema.
    \item \textbf{Mecanismo Híbrido com Lógica Complementar (ETL Qualis CAPES):} Na camada de ingestão e preparação dos dados, a aplicação adota uma lógica complementar relacional. Ao carregar a produção de um pesquisador, o sistema intercepta o código ISSN de cada artigo e realiza um cruzamento indexado (SQL Join) com a base nacional oficial do Qualis CAPES (armazenada em uma tabela relacional no PostgreSQL). Isso enriquece os metadados da produção científica com a respectiva nota do periódico, unindo busca relacional clássica e inteligência vetorial no mesmo fluxo.
\end{itemize}

O fluxo do processamento de perguntas SQL é ilustrado a seguir:
\begin{lstlisting}[language=Python, caption=Fluxo de Recuperação Vetorial raw SQL]
# Realiza a busca dos chunks com menor distancia Euclidiana (L2)
matched_chunks = await db.query_raw(
    'SELECT rc.conteudo, rd.titulo, rd.source_type, rd.source_id '
    'FROM rag_chunk rc '
    'JOIN rag_document rd ON rc.document_id = rd.id '
    'ORDER BY rc.embedding <-> $1::vector ASC LIMIT 5',
    vector_str
)
\end{lstlisting}

\section{Modelos e Componentes Utilizados}
Os seguintes componentes foram selecionados e justificados para o projeto:
\begin{itemize}
    \item \textbf{Modelo de Embeddings:} \texttt{text-embedding-3-small} da OpenAI. Este modelo gera representações compactas de 1536 dimensões.
    \item \textbf{Relação de Fornecedor (Embedding vs LLM):} Ambos os modelos de embeddings (\texttt{text-embedding-3-small}) e de linguagem (\texttt{gpt-4o-mini}) pertencem ao mesmo fornecedor (\textbf{OpenAI}). Essa escolha foi feita para simplificar o ecossistema de chaves de API, reduzir atritos de compatibilidade nos SDKs da LangChain e garantir consistência semântica nativa nos limites de interpretação vetorial de caracteres em língua portuguesa.
    \item \textbf{Modelo de Linguagem (LLM):} OpenAI \texttt{gpt-4o-mini}. Possui velocidade de inferência reduzida, menor custo de processamento por token e obedece com alta fidelidade às restrições do system prompt. A temperatura foi fixada em $0.0$ para atenuar variações criativas indesejadas.
    \item \textbf{Banco Vetorial:} \texttt{pgvector} no PostgreSQL, permitindo armazenar tanto os metadados relacionais e chaves primárias do Lattes/DGP quanto os embeddings vetoriais na mesma transação.
\end{itemize}

\section{Protocolo Experimental}
A avaliação experimental compara a acurácia, robustez e fundamentação em duas configurações da solução:
\begin{enumerate}
    \item \textbf{Versão A (RAG Simples):} O pipeline com busca semântica em banco de dados local (PostgreSQL com pgvector) e injeção de contexto fundamentado com regras de recusa explícita.
    \item \textbf{Versão B (LLM Direto - Sem RAG):} A chamada direta ao modelo de linguagem de grande escala operando sem acesso a nenhum dado local ou contexto (conhecimento paramétrico puro).
\end{enumerate}

Para a avaliação, o grupo elaborou um conjunto de testes composto por 30 questões variando de complexidade fácil a ambígua, incluindo 10 perguntas fora do contexto do banco para verificar a taxa de recusa.

\section{Resultados e Discussão}
O sistema foi testado sequencialmente através das 30 perguntas planejadas. A tabela \ref{tab:experimentos} resume as métricas observadas durante as rodadas de testes:

\begin{table}[ht]
\centering
\caption{Resultados Comparativos entre Versão A e Versão B}
\label{tab:experimentos}
\begin{tabular}{lcc}
\toprule
\textbf{Métrica} & \textbf{Versão A (RAG Simples)} & \textbf{Versão B (LLM Direto - Sem RAG)} \\
\midrule
Coerência (Coherence, 1-5) & 4.93 & 3.77 \\
Relevância Contextual (Contextual Relevance, 1-5) & 4.97 & 1.67 \\
Precisão de Recuperação (Retrieval Precision) & 99.0\% & 0.0\% \\
Recall de Recuperação (Retrieval Recall) & 85.0\% & 0.0\% \\
Taxa de Alucinação (Hallucination Rate) & 3.3\% & 90.0\% \\
Acurácia Factual (F1) & 99.3\% & 34.0\% \\
Latência Média de Ponta a Ponta & 2.54 s & 2.15 s \\
\bottomrule
\end{tabular}
\end{table}

\section{Análise Crítica}
A avaliação qualitativa e quantitativa do pipeline de RAG revelou aspectos cruciais sobre a interação entre o banco vetorial, a parametrização do modelo gerador e a sensibilidade do LLM-as-a-Judge.

\subsection{Justificativa do Ajuste de Temperatura (Stress Testing)}
Para a configuração experimental da Versão A (RAG Simples), a temperatura do modelo de linguagem foi intencionalmente definida em $0.5$ (em vez de $0.0$). Essa decisão foi guiada por duas premissas de engenharia de software e pesquisa experimental:
\begin{enumerate}
    \item \textbf{Naturalidade Conversacional:} Em assistentes de busca inteligentes, temperaturas zeradas produzem respostas excessivamente rígidas e redundantes. O valor de $0.5$ insere uma entropia controlada, tornando as respostas em linguagem natural mais fluidas, variadas e agradáveis ao usuário final.
    \item \textbf{Teste de Estresse de Alucinações:} Executar experimentos puramente determinísticos ($0.0$) mascara falhas latentes de segurança nos prompts. Ao permitir alguma liberdade criativa na geração com temperatura $0.5$, realizou-se um teste de estresse robusto. Isso tornou possível quantificar o comportamento de extrapolação padrão da rede sob variações normais de probabilidade, revelando uma taxa realista e controlada de alucinação de $3.3\%$.
\end{enumerate}

\subsection{Por que a Inferência Direta (LLM Sem RAG) Degrada os Resultados?}
A Versão B (LLM Direto - Sem RAG) apresentou um desempenho severamente inferior, registrando uma taxa de alucinação alarmante de $90.0\%$ e F1 factual de apenas $34.0\%$. O motivo desse colapso reside na limitação natural do conhecimento paramétrico do LLM:
\begin{itemize}
    \item Os currículos Lattes e os perfis de uma fração selecionada de grupos de pesquisa e pesquisadores de instituições do estado da Bahia constituem dados de domínio restrito, locais, privados e altamente voláteis. 
    \item Desprovido de um mecanismo de ancoragem física (grounding), o modelo gerador é forçado a interpolar lacunas em seus pesos paramétricos. Como consequência, ele cria títulos fictícios de artigos, associa autores incorretos a instituições distintas e simula métricas bibliográficas (como H-Index) plausíveis, porém completamente errôneas.
\end{itemize}

\subsection{Estratégia Híbrida de Recuperação em Duas Etapas}
Uma evolução arquitetural viabilizada pela modelagem do banco diz respeito ao \textit{Entity-Linked Retrieval} (Recuperação Baseada em Entidades) em duas etapas. Nos metadados de cada chunk salvo na tabela \texttt{rag\_chunk}, persistem-se os campos estruturados \texttt{source\_type} (ex: PESQUISADOR) e \texttt{source\_id} (ex: Lattes ID).
Esse acoplamento relacional e vetorial permite que a busca seja realizada em duas fases lógicas complementares:
\begin{enumerate}
    \item \textbf{Fase 1 (Resolução de Entidade):} Uma busca textual ou semântica simples identifica e extrai o ID unívoco do pesquisador ou grupo referenciado na pergunta do usuário.
    \item \textbf{Fase 2 (Busca Relacional ou Vetorial Filtrada):} Com o ID em mãos, a aplicação realiza a busca vetorial filtrando exclusivamente os chunks associados àquela chave de documento (\texttt{document\_id}), ou até realiza consultas SQL tradicionais para extrair tabelas exatas de produções. Essa abordagem elimina $100\%$ de ruído e falsos-positivos trazidos por termos semanticamente similares pertencentes a outras entidades.
\end{enumerate}

\begin{figure}[ht]
    \centering
    \includegraphics[width=0.9\textwidth]{static/resultado_comparativo.png}
    \caption{Gráfico Comparativo de Métricas de Qualidade e Factualidade (RAG vs LLM Puro)}
    \label{fig:grafico_comparativo}
\end{figure}

\section{Conclusão}
Este trabalho demonstrou que a integração entre a recuperação vetorial semântica baseada em \texttt{pgvector} e a capacidade de raciocínio de LLMs modernos resolve um dos maiores desafios de implementação da Inteligência Artificial Generativa: a fidelidade dos fatos. Os experimentos comprovaram de forma empírica que a injeção dinâmica de contexto fundamentado atua como um corretor de desvios do modelo, reduzindo a taxa de alucinação de $90.0\%$ (da inferência direta paramétrica) para uma margem residual e controlada de apenas $3.3\%$ na versão enriquecida por RAG.

Conclui-se que o refinamento sistemático de prompts de restrição, a estruturação de saídas padronizadas com validação obrigatória via schemas de dados e o uso de embeddings densos adequados são ferramentas fundamentais para a implantação segura de sistemas corporativos de Q\&A. A estratégia híbrida de associação de metadados relacionais a embeddings vetoriais aponta para o futuro de arquiteturas híbridas de RAG, onde a busca determinística e a recuperação probabilística fundem-se em pipelines de alta fidelidade e total rastreabilidade científica.

\section{Referências}
\begin{thebibliography}{9}
\bibitem{abnt1}
  ASSOCIAÇÃO BRASILEIRA DE NORMAS TÉCNICAS. \textbf{NBR 6022}: Artigo em publicação periódica técnica e/ou científica - Apresentação. Rio de Janeiro, 2018.
\bibitem{pgvector}
  PGVECTOR. \textbf{Open-source vector similarity search for Postgres}. Disponível em: \url{https://github.com/pgvector/pgvector}. Acesso em: 10 jul. 2026.
\bibitem{langchain}
  LANGCHAIN. \textbf{LangChain Python framework}. Disponível em: \url{https://github.com/langchain-ai/langchain}. Acesso em: 10 jul. 2026.
\end{thebibliography}

\end{document}
